Қалалық көлік басқару жүйелері (UTMS) сенсорлық және оқиға деректеріне сүйенеді. Бұл деректер болжам жасау немесе аномалия анықтау алгоритмдеріне дейін жиі толық емес, қайталанған немесе схема бойынша сәйкессіз болады. Бұл диссертацияда \textbf{Flowmatic} интеллектуалды платформасы әзірленеді, іске асырылады және бағаланады: ол деректерді дайындау мен inference процестерін автоматтандырады, сапаны бағалайды, дайындалған ағындарды модальдық нейрондық модельдерге бағыттайды және қайталанатын оқыту артефактілерін жариялайды.

Flowmatic контейнерленген веб-стекте алты өлшемді деректер сапасы индексін (DQI), автоматты тазалауды және экспортты қамтиды. Платформаны пакеттік тексеруде 30\,000 жолдық Астана деректер жиынтығы 406\,мс ішінде өңделеді; 20\,\%-ға дейінгі бақыланатын бүлінулер кезінде кешенді DQI 0{,}045-қа дейін қалпына келеді, ал жарамсыз жолдар жойылады. Smart City жұмыс ортасы симуляцияланған көлік пен ауа-райы дереккөздерін, құжатталған тізілім бойынша Manual/Auto маршруттауды және medallion форматындағы деректер көліне экспортты біріктіреді.

Жеті негізгі нейрондық модель ортақ уақыттық hold-out протоколымен (70/15/15 бөлу, үш seed) оқытылады. Ауырлық деңгейі классификаторы macro-F1 $= 0{,}911 \pm 0{,}017$ көрсеткішіне жетеді және сол hold-out тест бөлімінде sklearn негізіндегі базалық модельдерден асады; белгілер ереже негізіндегі синтетикалық атрибуттар. TranAD енгізілген терезелік аномалияларда AUROC $= 0{,}84$ көрсетеді. Дайындау абляциясы бүлінген ұяшықтар қалпына келтірілгенде 20\,\% бүлінуден кейін PatchTST density RMSE $0{,}90$-дан $0{,}58$-ға дейін төмендейді. 140 симулятор оқиғасында ереже-негізделген Auto маршруттау құжатталған эталондық сценарийлер бойынша 100\,\% сценарийлік саясатқа сәйкестік көрсетеді; бұл тәуелсіз эксперттік маршруттау дәлдігі емес.

Артефактілер Hugging Face платформасында \flowmaticrepo{} репозиторийлер отбасында жарияланған. Пакеттік API тексеруі тек Астана экспортын пайдаланады; нақты ведомстволарда орналастыру және сыртқы маршруттау белгілері болашақ жұмыс болып табылады.

\abstractkeywords{Тірек сөздер:}{деректерді дайындау; интеллектуалды көлік жүйелері; Smart City; Flowmatic; модель маршруттау; уақыт қатарлары; қайталанушылық}
